\section{Early decisions from componentwise Euler continuations}
\label{sec:euler}

\subsection{Why a continuation is needed}
Saturation alone permits an exact decision at closure. To stop earlier, we
must prove that enough of the current virtual state survives all remaining
crossings. A large current multiplicity does not do that. An exact Euler
calculation of each \emph{completed independent summand} supplies a safe
lower bound, without computing its full homology.

Let $F_i$ denote extension through the unprocessed suffix, followed by
closure. It is the same additive, homotopy-preserving operation used by the
ordinary scan. For a compressed state after $i$ crossings, define
\begin{equation}\label{eq:eulerbound}
 B_i=\sum_a p_a(1)\,|\chi(F_i(T_a))|.
\end{equation}
Absolute values are taken separately for the different templates before
combining their multiplicities. Using $p_a(-1)$ would allow copies in
different homological shifts to cancel and would lose this information.

\begin{theorem}[Monotone lower bound for final rank]\label{thm:euler}
Along a complete exact component-sharing run,
\[
 B_i\le r_2(K),\qquad B_i\le B_{i+1},\qquad B_n=r_2(K).
\]
Thus $B_i\ge3$ certifies \KNOTTED before closure. The cap
$\min(3,B_i)$ can be computed from saturated multiplicities, provided the
signed Euler characteristics themselves are computed exactly.
\end{theorem}
\begin{proof}
The direct-sum representation and preservation of homotopy equivalence give
\[
 r_2(K)=\sum_a p_a(1)\dim_{\F}H(F_i(T_a)).
\]
For any finite-dimensional complex $C$, the Euler identity and triangle
inequality give $|\chi(C)|\le\dim H(C)$. Applying this separately to the
completed templates proves the lower bound.

Suppose extension and reduction of one template produces
$E_i(T_a)\simeq\bigoplus_u U_u[e_u]$. Compatibility of continuation gives
\[
 \chi(F_i(T_a))=\sum_u(-1)^{e_u}\chi(F_{i+1}(U_u)).
\]
The triangle inequality, multiplied by $p_a(1)$ and summed over parents,
proves monotonicity. Merging equal normalized templates adds nonnegative
multiplicities and does not change the sum of absolute values. At final
closure, the fully eliminated complex is a direct sum of isolated objects;
every normalized template has Euler characteristic one. Hence $B_n$ is the
actual total rank.

Finally, multiplication by $|\chi|$ and summation are nonnegative
operations. Lemma~\ref{lem:semiring} shows that saturated multiplicities
suffice for the cap of~\eqref{eq:eulerbound}. The threshold follows
from~\eqref{eq:rankcriterion}.
\end{proof}

This is an obstruction to triviality. A value $B_i=2$ at an incomplete stage
is inconclusive. Some completed summands have nonzero homology and Euler
characteristic zero, so even genuine splitting need not produce an early
certificate.

\subsection{An exact suffix dynamic program}
For a matching $m$ on the frontier after crossing $i$, set
$E_i(m)=\chi(F_i(m))$. The final empty matching has
$E_n(\varnothing)=1$. This is the monoidal empty object; the separate API
convention that an empty \emph{input PD} is one unknot is handled separately.

Suppose the next crossing gives matching $m_0$ and $c_0$ closed circles in
smoothing zero, and matching $m_1$ and $c_1$ closed circles in smoothing one.
The two cube terms differ in degree by one, while each delooped circle has
two generators in that same homological degree. Therefore
\begin{equation}\label{eq:eulerrecurrence}
 E_i(m)=2^{c_0}E_{i+1}(m_0)-2^{c_1}E_{i+1}(m_1).
\end{equation}
For a normalized template,
\begin{equation}\label{eq:templateeuler}
 \chi(F_i(T))=\sum_{v\in T}(-1)^{h_v}E_i(m_v).
\end{equation}
The differential of $T$ is unnecessary for this particular calculation,
because Euler characteristic depends only on chain dimensions.

The implementation memoizes pairs consisting of the suffix index and the
full matching. It evaluates the dependency DAG using an explicit stack, so
a long narrow braid does not hit Python's recursion limit. It deliberately
uses~\eqref{eq:eulerrecurrence} rather than a writhe shortcut. The recurrence
follows the scanner's exact raw degree convention and works for arbitrary
valid scan permutations without additional orientation formulas.

\begin{proposition}[Cost of the Euler continuation]\label{prop:eulercost}
If all frontiers have size at most $w$, the complete memo table has at most
$(n+1)M(w)$ states. Each value has $O(n)$ bits. With ordinary finite-state
dynamic programming, all requested continuation values cost
$n^{O(1)}2^{O(w\log(w+1))}$ bit operations. Combining them with the physical
templates adds a polynomial cost in their stored size.
\end{proposition}
\begin{proof}
The boundary label set at each stage is fixed by the scan order, independently
of smoothing choices. There are at most $M(w)$ complete matchings on it.
Each state has two outgoing dependencies. From~\eqref{eq:eulerrecurrence}
and $c_0,c_1\le2$, induction gives $|E_i(m)|\le8^{n-i}$, so the bit bound
is linear. The two gluing operations are polynomial in $w$; memoized integer
addition and multiplication by powers of two have polynomial bit cost.
Use $M(w)=2^{O(w\log(w+1))}$. Full key equality instead of constant-time
hashing can be charged by a further polynomial factor in the state count,
which gives the same displayed exponential form.
\end{proof}

The coefficient values in~\eqref{eq:eulerrecurrence} must not be capped.
Subtraction may cancel two large terms to a small integer. Only after an
exact signed sum and an absolute value may the nonnegative lower bound be
saturated. This is tested explicitly, including mirror diagrams and
intermediate values larger than the final Euler magnitude.

The optional implementation uses a state budget, defaulting to 4096. If the
budget is exhausted, it disables only this inference and continues the
complete saturated scanner. A time or physical-object limit remains a limit
on the whole computation and becomes \UNKNOWN. The early result exposes
\code{rank\_lower\_bound\_capped}, not \code{rank} or
\code{rank\_capped}; the latter would misleadingly suggest a completed
homology calculation. The suffix-state budget includes all registered
states over the lifetime of the inference engine.

\subsection{Scope of the early-decision result}
The theorem supplies a new proof obligation that can be checked without
assuming that current objects survive: calculate the Euler characteristic
\emph{after the actual remaining continuation}. The bound is stronger than
the absolute Euler characteristic of the whole complex because opposite
Euler contributions from independent components are not cancelled.

It is still sensitive to the differential-component decomposition reached
by the scan. If there is one component with multiplicity one, its completed
Euler magnitude is just the ordinary knot Euler magnitude. The implementation
therefore waits for an actual split before attempting the optional dynamic
program. This is an optimization of a sound condition, not a completeness
assumption. An unbounded suffix calculation fits within the conservative
finite-type bound of Section~\ref{sec:complexity}; the budgeted version gives
better control of practical overhead.
